% ROZDZIAŁ 28 -- wzory miarowo-geometryczne, nie powtórzenie całek po podrozmaitościach.
\chapter{Wzór na pole powierzchni, krotności i koobszar}
\label{ch:area-coarea-28}

\section{Od wzoru parametrycznego do wzoru miarowego}
W skrypcie \emph{Geometria i topologia różniczkowa}
udowodniono wzór na całkę po kawałku gładkiej
podrozmaitości opisanej parametryzacją, która
pokrywa ten kawałek dokładnie raz. Tutaj rozważamy
\emph{odwzorowania nieiniektywne} oraz ich włókna.
Nowością są krotności obrazów, całkowanie po
przeciwobrazach punktów i formuły dla
odwzorowań Lipschitzowskich.

W tym rozdziale $\lambda_k$ oznacza miarę Lebesgue'a
na $\mathbb R^k$, a $\mathcal H^k$ -- miarę
Hausdorffa z Rozdziału~\ref{ch:hausdorff-26}.
Stosujemy konwencję z
Uwagi~\ref{rem:hausdorff-convention-26}: dla wymiaru
całkowitego $k\ge1$ wybieramy $c_k=\kappa_k^{-1}$,
tak aby $\mathcal H^k|_{\mathbb R^k}=\lambda_k$;
$\mathcal H^0$ jest miarą liczącą.
Dzięki temu poniższe wzory nie zakładają
nieudowodnionej nierówności izodiametrycznej.
Przy innej stałej normalizującej ich prawe strony
należy odpowiednio przeskalować.

\begin{definicja}[Jakobian wymiaru $k$]
\label{def:k-jacobian-28}
Niech $1\le k\le N$. Dla dowolnego odwzorowania
liniowego $A:\mathbb R^k\to\mathbb R^N$ definiujemy
\[
 J_k(A)=\sqrt{\det(A^{\mathsf T}A)}.
\]
Jeżeli $U\subseteq\mathbb R^k$ jest otwarty,
a $F:U\to\mathbb R^N$ jest różniczkowalne
w punkcie $x$, to kładziemy $J_kF(x)=J_k(DF(x))$.
W szczególności dotyczy to wszystkich punktów odwzorowania
$C^1$ oraz prawie wszystkich punktów odwzorowania
lokalnie Lipschitzowskiego. Jakobian jest zerowy,
gdy $\operatorname{rank}DF(x)<k$; dla $k=N$
jest równy $|\det DF(x)|$.
\end{definicja}

\begin{definicja}[Krotność odwzorowania]
\label{def:multiplicity-28}
Dla dowolnego odwzorowania $F:X\to Y$
i zbioru $E\subseteq X$ definiujemy
\[
 N(F,E,y)=\#\{x\in E:F(x)=y\}\in\mathbb N_0\cup\{+\infty\}.
\]
Symbol $\#$ oznacza moc zbioru, a nie operację
podnoszenia indeksów za pomocą metryki.
Pojawia się tu również wartość $+\infty$,
gdy włókno ma nieskończenie wiele punktów.
\end{definicja}

\section{Miara Hausdorffa gładkiego wykresu}
\begin{lemat}[Małe zniekształcenie wykresu]
\label{lem:graph-hausdorff-28}
Niech $V\subseteq\mathbb R^k$ będzie otwarty,
a $h:V\to\mathbb R^{N-k}$ funkcją $C^1$.
Niech $\Gamma(z)=(z,h(z))$ oraz
$\nu(A)=\mathcal H^k(\Gamma(A))$
dla borelowskich $A\subseteq V$. Wtedy $\nu$
jest miarą Radona na $V$ i dla każdego
$z_0\in V$ zachodzi
\[
 \lim_{r\downarrow0}
 \frac{\nu(B(z_0,r))}{\lambda_k(B(z_0,r))}
 =\sqrt{\det\bigl(I_k+(Dh(z_0))^{\mathsf T}Dh(z_0)\bigr)}.
\]
\end{lemat}
\begin{proof}
Na każdej kuli domkniętej $K\subset V$ wykres $\Gamma$
jest Lipschitzowski, ponieważ $Dh$ jest ograniczona
na $K$. Zatem, na mocy
Twierdzenia~\ref{thm:hausdorff-lipschitz-26},
$\mathcal H^k(\Gamma(K))<\infty$, gdyż
$\mathcal H^k(K)=\lambda_k(K)$ jest skończona.
Z drugiej strony rzut ortogonalny $\pi$ na
pierwszych $k$ współrzędnych jest odwrotnością
$\Gamma$ na obrazie i jest $1$-Lipschitzowski.
Ponieważ $\Gamma$ jest homeomorfizmem $V$ na
$\Gamma(V)$, przeciwobrazy zbiorów borelowskich
są borelowskie, a $\nu$ jest miarą borelowską
lokalnie skończoną. Stosujemy twierdzenie o
regularności miar lokalnie skończonych na lokalnie
zwartych ośrodkowych przestrzeniach metrycznych
z Rozdziału~\ref{ch:radon-measures-20}.
Otrzymujemy, że $\nu$ jest miarą Radona,
bez zakładania z góry regularności miary
Hausdorffa na całym wykresie.

Ustalmy $z_0\in V$ i liniowe odwzorowanie
$L(z)=(z,h(z_0)+Dh(z_0)(z-z_0))$.
Na małej kuli $B(z_0,r)$ z ciągłości $Dh$ wynika
$|[h(u)-h(v)]-Dh(z_0)(u-v)|\le\varepsilon_r|u-v|$,
gdzie $\varepsilon_r\to0$.
Ponieważ $|L(u)-L(v)|\ge|u-v|$,
odwzorowanie $\Gamma\circ L^{-1}$ między
odpowiednimi fragmentami wykresów jest
dwulipschitzowskie ze stałymi należącymi
do przedziału $[1-\varepsilon_r,1+\varepsilon_r]$
po ewentualnym zwiększeniu $\varepsilon_r$.
Z oszacowania miary Hausdorffa przez stałą
Lipschitza otrzymujemy
\[
 (1-\varepsilon_r)^k\mathcal H^k(L(B(z_0,r)))
 \le\nu(B(z_0,r))\le
 (1+\varepsilon_r)^k\mathcal H^k(L(B(z_0,r))).
\]
Liniowy obraz przestrzeni $\mathbb R^k$ jest
$k$-płaszczyzną euklidesową. Diagonalizacja
macierzy Grama pokazuje, że
$\mathcal H^k(L(A))=J_kL\,\lambda_k(A)$
dla każdego borelowskiego $A$, gdzie
$J_kL=\sqrt{\det(I_k+Dh(z_0)^{\mathsf T}Dh(z_0))}$.
Dzielimy przez $\lambda_k(B(z_0,r))$ i
przechodzimy do granicy.
\end{proof}

\begin{twierdzenie}[Miara na gładkim wykresie]
\label{thm:graph-area-28}
Przy założeniach Lematu~\ref{lem:graph-hausdorff-28}
dla każdej nieujemnej funkcji borelowskiej
$\psi$ na $\Gamma(V)$ mamy
\[
 \int_{\Gamma(V)}\psi(y)\,\mathrm d\mathcal H^k(y)
 =\int_V\psi(\Gamma(z))
 \sqrt{\det(I_k+Dh(z)^{\mathsf T}Dh(z))}
 \,\mathrm d\lambda_k(z).
\]
To samo zachodzi dla funkcji całkowalnych.
\end{twierdzenie}
\begin{proof}
Miara $\nu$ jest lokalnie skończona i
$\nu\ll\lambda_k$: obraz zbioru o zerowej
mierze Lebesgue'a ma zerową miarę Hausdorffa,
co wynika z lokalnej Lipschitzowskości $\Gamma$.
Istnieje zatem gęstość $a=\mathrm d\nu/\mathrm d\lambda_k$.
Twierdzenie Lebesgue'a o różniczkowaniu miar
(Rozdział~\ref{ch:differentiation-25})
identyfikuje $a(z)$ prawie wszędzie z granicą
ilorazów mas małych kul. Poprzedni lemat
pokazuje, że granica ta jest równa $J_k\Gamma(z)$.
Otrzymujemy więc równość miar
$\nu=J_k\Gamma\,\lambda_k$, a następnie
równość całek.
\end{proof}

\section{Wzór na pole powierzchni dla immersji}
\begin{twierdzenie}[Wzór z krotnością dla immersji $C^1$]
\label{thm:immersion-area-28}
Niech $1\le k\le N$, $U\subseteq\mathbb R^k$
będzie otwarty, a $F:U\to\mathbb R^N$
immersją klasy $C^1$, tj.
$\operatorname{rank}DF(x)=k$ dla każdego $x\in U$.
Niech $E\subseteq U$ będzie borelowski,
a $f:E\to[0,+\infty]$ borelowska.
Wówczas funkcja
$y\mapsto\sum_{x\in E\cap F^{-1}(\{y\})}f(x)$
jest $\mathcal H^k$-mierzalna i
\begin{equation}\label{eq:area-immersion-28}
 \int_E f(x)J_kF(x)\,\mathrm d\lambda_k(x)
 =\int_{\mathbb R^N}\left(
 \sum_{x\in E\cap F^{-1}(\{y\})}f(x)\right)
 \,\mathrm d\mathcal H^k(y).
\end{equation}
Przy $f\equiv1$ na $E$ (czyli $f=\chi_E$ na $U$)
prawa strona redukuje się do
\[
 \int_{\mathbb R^N} N(F,E,y)\,\mathrm d\mathcal H^k(y).
\]
\end{twierdzenie}
\begin{proof}
Twierdzenie o stałej randze zapewnia, że każdy
$x\in U$ ma otoczenie $U_x$, na którym $F$
jest osadzeniem na gładką podrozmaitość.
Przeliczalność bazy pozwala wybrać przeliczalne
pokrycie takimi otoczeniami $U_j$. Tworzymy
rozłączny podział borelowski
$B_1=U_1$, $B_j=U_j\setminus\bigcup_{i<j}U_i$.

Na każdym $U_j$ obraz $F(U_j)$ jest gładką
podrozmaitością w pewnym otoczeniu, a po
lokalnym przekształceniu ortogonalnym i zmianie
współrzędnych można go przedstawić jako wykres.
Twierdzenie~\ref{thm:graph-area-28}, zastosowane
w tych współrzędnych, oraz klasyczny wzór
podstawienia na dziedzinie $\mathbb R^k$
dają dla borelowskich $A\subseteq B_j$
\[
 \int_A f(x)J_kF(x)\,\mathrm d\lambda_k(x)
 =\int_{F(A)}f((F|_{U_j})^{-1}(y))
 \,\mathrm d\mathcal H^k(y).
\]
Różniczka zmiany parametrów wprowadza dokładnie
ten czynnik, który przekształca macierz Grama
$F$ w macierz Grama wykresu; jest to tożsamość
$J_k(F\circ\phi)=J_kF(\phi)|\det D\phi|$
dla lokalnego dyfeomorfizmu $\phi$.

Każdy $F(B_j\cap E)$ jest zbiorem borelowskim
w lokalnie osadzonym wykresie, a zatem również
w $\mathbb R^N$. Odpowiednie funkcje na obrazie
są borelowskie. Sumujemy równości po $j$;
po lewej używamy rozłączności $B_j$, a po prawej
twierdzenia Tonellego. Dla każdego $y$ suma
wkładów z kolejnych $B_j$ jest dokładnie sumą
po całym włóknie $E\cap F^{-1}(\{y\})$.
\end{proof}

\begin{przyklad}[Dlaczego trzeba liczyć krotności?]
Niech $F:(0,4\pi)\to\mathbb R^2$,
$F(t)=(\cos t,\sin t)$. Wówczas $J_1F=1$,
więc lewa strona \eqref{eq:area-immersion-28}
dla $f\equiv1$ wynosi $4\pi$.
Prawie każdy punkt okręgu ma dwa przeciwobrazy
w $(0,4\pi)$, zatem prawa strona wynosi
$2\mathcal H^1(S^1)=4\pi$.
Bez krotności otrzymalibyśmy wynik dwa razy za mały.
\end{przyklad}

\section{Jakobian normalny i całkowanie po włóknach}
\begin{definicja}[Jakobian normalny]
\label{def:normal-jacobian-28}
Niech $n\ge m\ge1$. Dla dowolnego odwzorowania
liniowego $A:\mathbb R^n\to\mathbb R^m$ definiujemy
\[
 \operatorname{Jac}_m(A)=\sqrt{\det(AA^{\mathsf T})}.
\]
Dla odwzorowania $F:U\to\mathbb R^m$ różniczkowalnego
w punkcie $x$ przyjmujemy
$\operatorname{Jac}_mF(x)=\operatorname{Jac}_m(DF(x))$.
Definicja obejmuje zatem odwzorowania $C^1$ i punkty
różniczkowalności odwzorowań Lipschitzowskich.
Jeżeli $\operatorname{rank}DF(x)<m$, to
$\operatorname{Jac}_mF(x)=0$; dla pełnej rangi
jest dodatni. Gdy $m=1$, jest równy $|\nabla F(x)|$.
\end{definicja}

\begin{lemat}[Tożsamość liniowo-algebraiczna koobszaru]
\label{lem:linear-coarea-28}
Niech $A:\mathbb R^n\to\mathbb R^m$ będzie
surjektywnym odwzorowaniem liniowym, a
$G:\mathbb R^m\times\mathbb R^{n-m}\to\mathbb R^n$
odwracalnym odwzorowaniem liniowym, dla którego
$AG(y,z)=y$. Wówczas
\[
 \sqrt{\det(AA^{\mathsf T})}\,|\det DG|
 =\sqrt{\det((D_zG)^{\mathsf T}D_zG)}.
\]
\end{lemat}
\begin{proof}
Wybieramy ortonormalne rozkłady
$\mathbb R^n=(\ker A)^\perp\oplus\ker A$
i zapisujemy macierze względem tych baz.
Ograniczenie $A_0=A|_{(\ker A)^\perp}$
jest izomorfizmem, a macierz $G$ ma blokową postać
$\begin{psmallmatrix} A_0^{-1}&0\\B&C\end{psmallmatrix}$,
gdzie $C$ jest odwracalną macierzą kwadratową
rozmiaru $n-m$. Stąd
$|\det DG|=|\det A_0|^{-1}|\det C|$.
Ponieważ $AA^{\mathsf T}=A_0A_0^{\mathsf T}$,
lewy pierwiastek jest równy $|\det A_0|$.
Kolumny $D_zG$ leżą w $\ker A$, a ich macierz
w wybranej ortonormalnej bazie jest $C$,
więc prawy pierwiastek wynosi $|\det C|$.
Dla $n=m$ przyjmujemy wyznacznik pustej
macierzy równy $1$ i dowód pozostaje poprawny.
\end{proof}

\begin{twierdzenie}[Wzór koobszaru dla submersji]
\label{thm:coarea-submersion-28}
Niech $n\ge m\ge1$, $U\subseteq\mathbb R^n$
będzie otwarty, a $F:U\to\mathbb R^m$
submersją klasy $C^1$, tj.
$\operatorname{rank}DF(x)=m$ dla każdego $x\in U$.
Dla każdej borelowskiej funkcji $g:U\to[0,+\infty]$
zachodzi
\begin{equation}\label{eq:coarea-smooth-28}
 \int_U g(x)\operatorname{Jac}_mF(x)\,\mathrm d\lambda_n(x)
 =\int_{\mathbb R^m}\left(
 \int_{F^{-1}(\{y\})}g(x)\,\mathrm d\mathcal H^{n-m}(x)
 \right)\mathrm d\lambda_m(y).
\end{equation}
Wewnętrzna całka dotyczy miary Hausdorffa
na włóknie, które jest gładką
$(n-m)$-wymiarową podrozmaitością (a dla
$n=m$ jest zliczaniem punktów). Jeżeli jedna ze stron
jest skończona dla $|g|$, równość zachodzi także
dla rzeczywistych i zespolonych funkcji całkowalnych.
\end{twierdzenie}
\begin{proof}
Gdy $n=m$, odwzorowanie $F$ jest lokalnym
dyfeomorfizmem. Dziedzinę dzielimy na przeliczalnie
wiele części borelowskich, na których $F$ jest
różnowartościowe. Zwykła zamiana zmiennych i
Tonelli dają tezę, ponieważ
$\mathcal H^0$ jest miarą liczącą. W dalszym
dowodzie przyjmujemy zatem $n>m$.

\emph{Krok 1: lokalne współrzędne włóknowe.}
Na mocy twierdzenia o funkcji uwikłanej, dla
każdego $x_0\in U$ istnieje otoczenie $V$ i
dyfeomorfizm $G:W\to V$, gdzie
$W\subseteq\mathbb R^m\times\mathbb R^{n-m}$,
taki że $F(G(y,z))=y$ dla $(y,z)\in W$.
Wybieramy $g$ o nośniku zawartym w takim $V$.

\emph{Krok 2: zamiana zmiennych i Fubini.}
Przez klasyczną zamianę zmiennych dostajemy
\[
 \int_V g(x)\operatorname{Jac}_mF(x)\,\mathrm d\lambda_n(x)
 =\int_W g(G(y,z))\operatorname{Jac}_mF(G(y,z))
 |\det DG(y,z)|\,\mathrm d\lambda_n(y,z).
\]
Ponieważ $DF(G)D_{(y,z)}G=(I_m\ \ 0)$,
Lemat~\ref{lem:linear-coarea-28} zastosowany
punktowo do różniczek daje
$\operatorname{Jac}_mF(G)|\det DG|=
J_{n-m}(z\mapsto G(y,z))$.
Przez Tonellego całka jest więc równa
\[
 \int_{\mathbb R^m}\left(\int_{W_y}
 g(G(y,z))J_{n-m}(z\mapsto G(y,z))
 \,\mathrm d\lambda_{n-m}(z)\right)\mathrm d\lambda_m(y),
\]
gdzie $W_y=\{z:(y,z)\in W\}$.
Odwzorowanie $z\mapsto G(y,z)$ jest gładkim
osadzeniem w $F^{-1}(\{y\})$. Stosujemy
Twierdzenie~\ref{thm:immersion-area-28} na włóknie
(bez krotności) i otrzymujemy prawą stronę
\eqref{eq:coarea-smooth-28} dla $g$ wspieranej w $V$.

\emph{Krok 3: sklejenie.}
Wybieramy przeliczalne pokrycie $U$ otoczeniami
$V_j$ jak w kroku 1 i rozłączny podział borelowski
$B_1=V_1$, $B_j=V_j\setminus\bigcup_{i<j}V_i$.
Zastępujemy $g$ przez $g\chi_{B_j}$, sumujemy
równości lokalne po $j$ i stosujemy Tonellego
po obu stronach. Dla $n=m$ wewnętrzne całkowanie
względem $\mathcal H^0$ oznacza sumowanie po
włóknach; argument z lokalnymi odwrotnościami
pozostaje poprawny.
\end{proof}

\begin{wniosek}[Różniczkowanie poziomic]
\label{cor:coarea-level-28}
Dla $U\subseteq\mathbb R^n$ otwartego i
funkcji $F\in C^1(U;\mathbb R)$
bez punktów krytycznych oraz dowolnej
nieujemnej funkcji borelowskiej $g$ zachodzi
\[
 \int_U g(x)|\nabla F(x)|\,\mathrm d\lambda_n(x)
 =\int_{\mathbb R}\left(\int_{F^{-1}(\{t\})}
 g(x)\,\mathrm d\mathcal H^{n-1}(x)\right)
 \,\mathrm d\lambda_1(t).
\]
\end{wniosek}

\section{Zastosowanie: rozkład miary obrazowej}
\begin{twierdzenie}[Gęstość obrazu przez submersję]
\label{thm:pushforward-coarea-28}
Niech $F:U\to\mathbb R^m$ spełnia założenia
Twierdzenia~\ref{thm:coarea-submersion-28},
a $w:U\to[0,+\infty]$ będzie borelowska.
Niech $\nu=w\lambda_n$ i $\eta=F_*\nu$.
Wówczas $\eta\ll\lambda_m$, a jej gęstość jest
równa
\[
 q(y)=\int_{F^{-1}(\{y\})}
 \frac{w(x)}{\operatorname{Jac}_mF(x)}
 \,\mathrm d\mathcal H^{n-m}(x)
\]
dla $\lambda_m$-prawie wszystkich $y$,
z możliwością przyjmowania wartości $+\infty$.
\end{twierdzenie}
\begin{proof}
Dla borelowskiego $A\subseteq\mathbb R^m$
stosujemy wzór koobszaru do funkcji
\[
 g(x)=\frac{\chi_{F^{-1}(A)}(x)w(x)}
 {\operatorname{Jac}_mF(x)}.
\]
Mianownik jest dodatni na całym $U$, więc
jest to funkcja borelowska nieujemna.
Wynik ma postać
$\eta(A)=\int_A q(y)\,\mathrm d\lambda_m(y)$.
To zarazem dowodzi absolutnej ciągłości
i identyfikuje pochodną Radona--Nikodyma.
\end{proof}

\begin{przyklad}[Poziomice promienia w pierścieniu]
Niech $U=\{x\in\mathbb R^2:a<|x|<b\}$,
gdzie $0<a<b$, i $F(x)=|x|^2$.
Wtedy $|\nabla F(x)|=2|x|$. Coarea daje
\[
 \int_U 2|x|\,\mathrm d\lambda_2(x)
 =\int_{a^2}^{b^2}\mathcal H^1(F^{-1}(\{t\}))
 \,\mathrm d\lambda_1(t)
 =\int_{a^2}^{b^2}2\pi\sqrt t\,\mathrm d\lambda_1(t)
 =\frac{4\pi}{3}(b^3-a^3).
\]
Jest to ilustracja całkowania po \emph{włóknach},
a nie kolejny sposób obliczania powierzchni sfery.
\end{przyklad}


\section{Twierdzenie Rademachera: skąd bierze się różniczka prawie wszędzie?}
\label{sec:rademacher-28}

Wzory geometryczne wymagają jakobianu, czyli różniczki. W przypadku
odwzorowań klasy $C^1$ istnieje ona wszędzie, natomiast odwzorowanie
Lipschitzowskie może nie być różniczkowalne w pewnych punktach.
Twierdzenie Rademachera dowodzi, że taki zbiór wyjątkowy ma zerową
miarę Lebesgue'a. Samo istnienie różniczki prawie wszędzie nie dowodzi
jednak ogólnych wzorów area i coarea: do ich uzasadnienia potrzebne są
jeszcze twierdzenia o pokrywaniu i mierzeniu obrazów zbiorów krytycznych.

\begin{definicja}[Różniczkowalność Frécheta]
\label{def:frechet-rademacher-28}
Niech $U\subseteq\mathbb R^n$ będzie otwarty,
$f:U\to\mathbb R^m$ i $x\in U$. Mówimy, że $f$ jest
\emph{różniczkowalna w sensie Frécheta} w $x$, jeżeli istnieje
liniowe $A:\mathbb R^n\to\mathbb R^m$ takie, że
\[
 \lim_{h\to0}\frac{|f(x+h)-f(x)-Ah|}{|h|}=0.
\]
Takie $A$ jest jednoznaczne i oznaczamy je przez $Df(x)$.
Jeżeli dla każdego $x\in K\Subset U$ istnieją $L_K$ oraz
otoczenie $V_K$ zbioru $K$, na którym
$|f(x)-f(y)|\le L_K|x-y|$, nazywamy $f$
\emph{lokalnie Lipschitzowskim}. Równoważnie każdy punkt
$U$ ma otoczenie, na którym obowiązuje taka nierówność
z pewną stałą.
\end{definicja}

\begin{lemat}[Jednowymiarowa różniczkowalność]
\label{lem:rademacher-line-28}
Niech $a<b$, a $f:[a,b]\to\mathbb R$ będzie funkcją
Lipschitzowską ze stałą $L$. Wtedy istnieje
$g\in L^\infty((a,b),\lambda_1)$ takie, że
$|g|\le L$ prawie wszędzie oraz
\[
 f(t)-f(a)=\int_{(a,t)}g(s)\,\mathrm d\lambda_1(s)
 \quad (t\in[a,b]).
\]
W konsekwencji $f'=g$ prawie wszędzie.
\end{lemat}
\begin{proof}
Ponieważ $f$ jest Lipschitzowska, jest absolutnie ciągła.
Dla przejrzystości odzyskamy gęstość bez powoływania się na
odwrotne podstawowe twierdzenie rachunku całkowego.
Połóżmy $p(t)=f(t)+Lt$. Dla $s<t$ mamy
$p(t)-p(s)\ge -L(t-s)+L(t-s)=0$, więc $p$ jest
niemalejąca i ciągła. Jej przyrosty na półotwartych
przedziałach określają dodatnią miarę borelowską
$\mu_p((s,t])=p(t)-p(s)$; istnienie tej miary wynika
z konstrukcji Lebesgue'a--Stieltjesa przez rozszerzenie
Carathéodory'ego (ciągłość $p$ zapewnia ciągłość
premiary z góry na przedziałach). Dla każdego przedziału
$[s,t]$ zachodzi
$0\le\mu_p((s,t])\le2L(t-s)$.
Twierdzenie o jednoznaczności miar na przedziałach daje
$\mu_p\le2L\lambda_1$. Twierdzenie Radona--Nikodyma
dostarcza więc $q\in L^\infty$ o wartościach w $[0,2L]$
takiego, że $\mu_p=q\lambda_1$. Zatem
$f(t)-f(a)=\int_{(a,t)}(q(s)-L)\,\mathrm d\lambda_1(s)$.
Dla $g=q-L$ zachodzi $|g|\le L$ prawie wszędzie;
Twierdzenie~\ref{thm:FTC-lebesgue-25} kończy dowód.
\end{proof}

\begin{lemat}[Słabe i kierunkowe pochodne funkcji Lipschitzowskiej]
\label{lem:rademacher-direction-28}
Niech $f:U\to\mathbb R$ będzie lokalnie Lipschitzowska.
Istnieją mierzalne funkcje $g_i\in L^\infty_{\mathrm{loc}}(U)$,
$i=1,\ldots,n$, takie, że dla $\varphi\in C_c^1(U)$
\[
 \int_U f(x)\partial_i\varphi(x)\,\mathrm d\lambda_n(x)
 =-\int_U g_i(x)\varphi(x)\,\mathrm d\lambda_n(x).
\]
Dla każdego ustalonego $v\in\mathbb R^n$ pochodna
kierunkowa $\partial_v f(x)=\lim_{t\to0}
(f(x+tv)-f(x))/t$ istnieje dla prawie wszystkich
$x\in U$ i jest tam równa $\sum_i v_i g_i(x)$.
\end{lemat}
\begin{proof}
Przypuśćmy najpierw, że $f$ jest Lipschitzowska na
prostopadłościanie $Q\Subset U$. Z Lematu
\ref{lem:rademacher-line-28} i Fubiniego otrzymujemy
istnienie $\partial_i f$ prawie wszędzie na $Q$ oraz
$|\partial_i f|\le L$. Jest to funkcja mierzalna,
ponieważ jest granicą mierzalnych ilorazów różnicowych
w punktach istnienia granicy. Niech $\varphi\in C_c^1(Q)$.
Na prawie każdej prostej równoległej do osi $i$ funkcja
$f$ jest absolutnie ciągła. Jednowymiarowe całkowanie
przez części i Fubini dają
$\int_Q f\partial_i\varphi\,\mathrm d\lambda_n
=-\int_Q (\partial_i f)\varphi\,\mathrm d\lambda_n$.
Wyczerpując $U$ przeliczalnie wieloma takimi
prostopadłościanami, otrzymujemy globalną funkcję
$g_i$ z lokalnym ograniczeniem.

Dla dowolnego ustalonego $v\ne0$ rozkładamy
$\mathbb R^n=v^\perp\oplus\mathbb Rv$.
Po odpowiedniej ortogonalnej zmianie współrzędnych
stosujemy jednowymiarowy lemat do przekrojów
w kierunku $v$ i twierdzenie Fubiniego.
Wynika stąd, że $\partial_v f$ istnieje prawie
wszędzie oraz że jej działanie dystrybucyjne jest
równe $\partial_v f=\sum_i v_i g_i$.
Jednoznaczność funkcji lokalnie całkowalnej
reprezentującej dystrybucję dowodzi ostatniej równości.
\end{proof}

\begin{twierdzenie}[Rademacher]
\label{thm:rademacher-28}
Niech $n,m\ge1$, $U\subseteq\mathbb R^n$ będzie
otwarty, a $f:U\to\mathbb R^m$ lokalnie Lipschitzowska.
Wtedy istnieje zbiór borelowski $N\subseteq U$ o
$\lambda_n(N)=0$ taki, że $f$ jest różniczkowalna
w sensie Frécheta w każdym punkcie $U\setminus N$.
Ponadto w każdym punkcie różniczkowalności
$\|Df(x)\|\le\operatorname{Lip}(f;V)$ dla
każdego otoczenia $V$ punktu $x$, na którym $f$
jest Lipschitzowska.
\end{twierdzenie}
\begin{proof}
Wystarczy przypadek $m=1$, ponieważ stosujemy go
do wszystkich współrzędnych funkcji wektorowej.
Niech $g_1,\ldots,g_n$ będą słabymi pochodnymi z
Lematu~\ref{lem:rademacher-direction-28}. Rozważmy
przeliczalny zbiór kierunków $D=\mathbb Q^n\setminus\{0\}$.
Dla każdego $v\in D$ na zbiorze pełnej miary
istnieje $\partial_v f(x)=\sum_i v_i g_i(x)$.
Przecięcie tych przeliczalnie wielu zbiorów nadal
ma pełną miarę. Ustalmy punkt $x$ tego przecięcia,
w którym wszystkie $g_i(x)$ są skończone, i liniowe
odwzorowanie $A_xh=\sum_i g_i(x)h_i$.

Pokażemy, że $f(x+h)=f(x)+A_xh+o(|h|)$.
Ustalmy $\varepsilon>0$. Z gęstości racjonalnych
kierunków na sferze można wybrać \emph{skończenie}
wiele wektorów jednostkowych
$q_1,\ldots,q_M\in\{v/|v|:v\in D\}$, tak aby
każdy $u\in S^{n-1}$ był w odległości mniejszej
od $\delta$ od któregoś $q_j$, gdzie
$\delta>0$ wybierzemy później. Dla $h=tu$, $t=|h|$,
wybieramy taki $q_j$. Jeżeli $L$ jest stałą
Lipschitza $f$ w otoczeniu $x$, to
\begin{align*}
 |f(x+h)-f(x)-A_xh|
 &\le |f(x+tu)-f(x+tq_j)|\\
 &\quad+|f(x+tq_j)-f(x)-tA_xq_j|
       +|tA_x(q_j-u)|\\
 &\le t(L+\|A_x\|)\delta+o_j(t).
\end{align*}
Ponieważ lista kierunków $q_j$ jest skończona,
po podzieleniu przez $t$ ilorazy $o_j(t)/t$
zbiegają do zera \emph{jednostajnie względem $j$}.
Wybieramy $\delta$ tak, aby
$(L+\|A_x\|)\delta<\varepsilon$, a potem
dostatecznie małe $t$. Wynika stąd
$\limsup_{h\to0}|f(x+h)-f(x)-A_xh|/|h|\le\varepsilon$.
Dowolność $\varepsilon$ daje różniczkowalność
Frécheta z $Df(x)=A_x$.

Pozostaje uzasadnić borelowskość zbioru wyjątkowego.
Dla $A\in\mathbb Q^{m\times n}$, $r>0$ i $x\in U$
połóżmy
\[
 R_{A,r}(x)=\sup_{\substack{h\in\mathbb Q^n,\;0<|h|<r\\
                                    x+h\in U}}
 \frac{|f(x+h)-f(x)-Ah|}{|h|},
\]
przy czym supremum pustego zbioru przyjmujemy równe
zeru. Dla ustalonych $A,r$ funkcja $R_{A,r}$ jest
borelowska: supremum liczymy po przeliczalnie wielu
borelowskich ilorazach na otwartych dziedzinach
$\{x\in U:x+h\in U\}$.
Punkt $x$ jest różniczkowalnościowy wtedy i tylko wtedy,
gdy dla każdego $q\ge1$ istnieją
$A\in\mathbb Q^{m\times n}$ i $j\ge1$ takie, że
$R_{A,1/j}(x)<1/q$. Istotnie, jeżeli $Df(x)$ istnieje,
przybliżamy go macierzą wymierną. W drugą stronę
wybrane macierze $A_q$ spełniają
$\|A_q-A_p\|\le1/q+1/p$:
wspólne małe przyrosty wymierne, po podzieleniu przez
ich długość i przejściu do granicy, oszacowują normę
różnicy tych macierzy. Stąd $A_q\to A$, a z nierówności
resztowej wynika $f(x+h)=f(x)+Ah+o(|h|)$, najpierw
dla wymiernych $h$, potem dla wszystkich $h$ przez
ciągłość $f$. To daje przeliczalny opis borelowski
zbioru różniczkowalności. Jego dopełnienie ma zerową
miarę na mocy poprzedniego argumentu.
Oszacowanie normy różniczki wynika z przejścia do
granicy w nierówności Lipschitza.
\end{proof}

\begin{uwaga}[Dlaczego pochodne cząstkowe nie wystarczają?]
Na przykład $f(x,y)=|x|$ ma pochodne cząstkowe
w punktach z $x\ne0$, ale nie ma różniczki na prostej
$x=0$; zbiór ten jest zerowy dla $\lambda_2$.
Ponadto istnienie wszystkich pochodnych kierunkowych
w \emph{jednym} punkcie nie zapewnia różniczkowalności
Frécheta: dla $f(x,y)=\sqrt{x^2+y^2}$ w $(0,0)$
istnieją jednostronne pochodne kierunkowe, które nie
tworzą odwzorowania liniowego. Istotna w dowodzie
Rademachera jest zgodność pochodnych kierunkowych
z jedną \emph{słabą} różniczką, a następnie
argument skończonej siatki kierunków.
\end{uwaga}

\begin{przyklad}[Wszystkie pochodne cząstkowe w punkcie, ale brak różniczki]
Niech $f:\mathbb R^2\to\mathbb R$ będzie zadana przez
$f(0,0)=0$ oraz $f(x,y)=xy/\sqrt{x^2+y^2}$
dla $(x,y)\ne(0,0)$. Poza początkiem ma ona
ograniczony gradient: jest jednorodna stopnia $1$,
a jej pochodne pierwszego rzędu są ciągłe na
okręgu jednostkowym i jednorodne stopnia $0$.
W związku z tym $f$ jest globalnie Lipschitzowska
(dla odcinka przechodzącego przez $0$ dzielimy go
na dwa odcinki). Obie pochodne cząstkowe w $(0,0)$
istnieją i są równe $0$, ponieważ $f(x,0)=f(0,y)=0$.
Jednak dla $h=(t,t)$, $t>0$, mamy
$f(t,t)=t/\sqrt2$, a więc
$|f(h)-f(0)-0\cdot h|/|h|=1/2$.
Nie istnieje zatem różniczka Frécheta w $(0,0)$.
Pokazuje to, dlaczego sam argument z istnienia
pochodnych cząstkowych prawie wszędzie jest
niewystarczający dla twierdzenia Rademachera.
\end{przyklad}

\begin{przyklad}[Różniczkowalność prawie wszędzie a brak gładkości]
Niech $F:\mathbb R^2\to\mathbb R^2$,
$F(x,y)=(x,y+|x|)$. Jest globalnie Lipschitzowskie,
choć nie jest klasy $C^1$. Dla $x>0$ i $x<0$
ma odpowiednio macierze $\begin{psmallmatrix}1&0\\1&1\end{psmallmatrix}$
i $\begin{psmallmatrix}1&0\\-1&1\end{psmallmatrix}$,
a na prostej $x=0$ nie ma różniczki. Mimo to
$|\det DF|=1$ prawie wszędzie. Wzór zamiany
zmiennych jest poprawny, ale jego uzasadnienie
w klasie Lipschitzowskiej nie może ograniczać się
do zwykłego twierdzenia o dyfeomorfizmach $C^1$.
\end{przyklad}

\begin{uwaga}[Związek z miarami wektorowymi]
Twierdzenie Rademachera wykorzystuje skończony wymiar
\emph{przestrzeni wartości}. Dla dowolnej przestrzeni
Banacha $X$ nie musi być prawdą, że każda funkcja
Lipschitzowska $[0,1]\to X$ ma pochodną prawie wszędzie.
Dokładną granicę wyznacza \emph{własność Radona--Nikodyma},
która zostanie zdefiniowana i scharakteryzowana w
Rozdziale~\ref{ch:vector-measures-30}.
\end{uwaga}

\section{Aproksymacja typu Luzina--Whitneya w klasie $C^1$}
Aby przejść od wersji gładkich do odwzorowań Lipschitzowskich,
nie wystarczy różniczkowalność prawie wszędzie. Potrzebujemy
możliwości zastąpienia odwzorowania na dużych zbiorach
\emph{dokładnie} funkcją klasy $C^1$, wraz z jej różniczką.
Poniżej dowodzimy odpowiedniego faktu, nie przyjmując
jako gotowego twierdzenia o przedłużaniu Whitneya.

\begin{lemat}[Przedłużenie zgodnych danych pierwszego rzędu]
\label{lem:whitney-c1-28}
Niech $K\subseteq\mathbb R^n$ będzie niepustym zbiorem zwartym,
$a:K\to\mathbb R^N$ i
$A:K\to\mathcal L(\mathbb R^n,\mathbb R^N)$ funkcjami ciągłymi.
Załóżmy, że
\begin{equation}\label{eq:whitney-condition-28}
 \lim_{r\downarrow0}\ \sup_{\substack{x,y\in K,\;0<|x-y|\le r}}
 \frac{|a(y)-a(x)-A(x)(y-x)|}{|y-x|}=0,
\end{equation}
przy czym supremum pustego zbioru jest równe zeru.
Istnieje funkcja $\widetilde a\in C^1(\mathbb R^n;\mathbb R^N)$
taka, że $\widetilde a|_K=a$ i
$D\widetilde a(x)=A(x)$ dla $x\in K$.
\end{lemat}
\begin{proof}
Wprowadzamy dwa moduły ciągłości: $\omega(t)$ jest
supremum z \eqref{eq:whitney-condition-28} z $r=t$,
a $\alpha(t)=\sup\{\|A(x)-A(y)\|:x,y\in K,
|x-y|\le t\}$. Obydwa dążą do zera dla $t\downarrow0$.

\emph{Krok 1: kostki otwartego dopełnienia.}
Zbiór $\mathbb R^n\setminus K$ rozkładamy na kostki
diadyczne $Q$ o rozłącznych wnętrzach, dobierając
kostki maksymalne (względem włączenia w rodzinie
kostek diadycznych) spośród tych, dla których
$20\sqrt n\,\ell(Q)\le\operatorname{dist}(Q,K)$.
Tutaj $\ell(Q)$ oznacza długość boku kostki.
Każdy punkt dopełnienia należy do jakiejś kostki:
dostatecznie mała kostka diadyczna zawierająca ten
punkt spełnia nierówność. Maksymalność zapewnia
również oszacowanie odwrotne
$\operatorname{dist}(Q,K)\le C_n\ell(Q)$:
gdyby odległość była dużo większa, kostka macierzysta
(o dwukrotnie większym boku) też spełniałaby warunek.
Zatem długość boku jest porównywalna z odległością
kostki od $K$. W szczególności kostki o przecinających
się nieznacznie powiększonych wnętrzach mają
porównywalne boki, a liczba takich kostek w otoczeniu
dowolnego punktu jest ograniczona przez stałą zależną
wyłącznie od $n$ (wynika to z rozłączności wnętrz
i porównania objętości).

Dla każdej kostki wybieramy $p_Q\in K$ minimalizujący
odległość od jej środka. Wybieramy $C^\infty$-funkcje
$\psi_Q\ge0$, równe $1$ na $Q$, o nośnikach zawartych
w kostkach współśrodkowych o boku $(9/8)\ell(Q)$,
i spełniające $\|D\psi_Q\|\le C_n/\ell(Q)$.
Można je otrzymać jako iloczyny przeskalowanych
jednowymiarowych funkcji odcinających.
Suma $\sum_Q\psi_Q$ jest lokalnie skończona i nie
mniejsza niż $1$ na $\mathbb R^n\setminus K$.
Funkcje $\varphi_Q=\psi_Q/(\sum_R\psi_R)$ tworzą
więc gładki podział jedności, dla którego
$\|D\varphi_Q\|\le C_n/\ell(Q)$.

\emph{Krok 2: sklejanie przybliżeń afinicznych.}
Dla $x\notin K$ definiujemy
\[
 P_Q(x)=a(p_Q)+A(p_Q)(x-p_Q),\qquad
 \widetilde a(x)=\sum_Q\varphi_Q(x)P_Q(x),
\]
a dla $x\in K$ przyjmujemy $\widetilde a(x)=a(x)$.
Niech $x\notin K$, $\delta=\operatorname{dist}(x,K)$,
a $p\in K$ będzie punktem najbliższym $x$.
Dla każdej kostki, której $\varphi_Q(x)\ne0$,
mamy $\ell(Q)\asymp\delta$ i $|p_Q-p|\le C_n\delta$.
Warunek \eqref{eq:whitney-condition-28} daje
\[
 |P_Q(x)-[a(p)+A(p)(x-p)]|
 \le C_n\delta\,[\omega(C_n\delta)+\alpha(C_n\delta)],
\]
zaś $\|A(p_Q)-A(p)\|\le\alpha(C_n\delta)$.
Ponieważ $\sum_QD\varphi_Q=0$, różniczkując
sumę na dopełnieniu, otrzymujemy
\begin{align*}
 |\widetilde a(x)-a(p)-A(p)(x-p)|
 &\le C_n\delta[\omega(C_n\delta)+\alpha(C_n\delta)],\\
 \|D\widetilde a(x)-A(p)\|
 &\le C_n[\omega(C_n\delta)+\alpha(C_n\delta)].
\end{align*}
Stałe mogą również zależeć od wyboru funkcji
odcinających, ale nie od $x$.

Gdy $x\to p_0\in K$ od zewnątrz, mamy $p\to p_0$
i $\delta\to0$; druga nierówność zapewnia
$D\widetilde a(x)\to A(p_0)$. Z pierwszej oraz
\eqref{eq:whitney-condition-28} wynika
$\widetilde a(x)-a(p_0)-A(p_0)(x-p_0)=o(|x-p_0|)$.
Ten sam warunek bezpośrednio daje taki rozwój,
gdy $x\in K$. Zatem $\widetilde a$ jest
różniczkowalna także na $K$, jej różniczka jest
tam równa $A$ i jest ciągła na całej przestrzeni.
\end{proof}

\begin{lemat}[Zwarte fragmenty klasy $C^1$]
\label{lem:lusin-whitney-28}
Niech $U\subseteq\mathbb R^n$ będzie otwarty,
$F:U\to\mathbb R^N$ lokalnie Lipschitzowski,
$Q\subseteq U$ zwarty i $\varepsilon>0$.
Istnieją zbiór zwarty $K\subseteq Q$ oraz
$F_K\in C^1(\mathbb R^n;\mathbb R^N)$ takie, że
\[
 \lambda_n(Q\setminus K)<\varepsilon,\qquad
 F_K(x)=F(x),\quad DF_K(x)=DF(x)\quad(x\in K).
\]
Różniczkę $DF$ rozumiemy tylko w punktach,
w których istnieje; wszystkie punkty $K$ są
takimi punktami. Ponadto $U$ jest sumą przeliczalnej
rodziny zbiorów zwartych $K_j$, z których każdy
ma opisane wyżej przedłużenie, a
$\lambda_n(U\setminus\bigcup_jK_j)=0$.
\end{lemat}
\begin{proof}
Niech $D\subseteq U$ będzie borelowskim zbiorem
punktów różniczkowalności $F$. Na mocy
Twierdzenia~\ref{thm:rademacher-28} dopełnienie
$D$ jest zerowe. Funkcja $x\mapsto DF(x)$ jest
borelowska na $D$, ponieważ jej współczynniki
są granicami ilorazów różnicowych w kierunkach
wektorów bazowych.

Dla $j,l\ge1$ oznaczmy przez $H_{j,l}$ zbiór
punktów $x\in D\cap Q$, dla których
\[
 |F(x+h)-F(x)-DF(x)h|\le j^{-1}|h|
\]
dla każdego $h\in\mathbb Q^n$ spełniającego
$0<|h|<l^{-1}$ i $x+h\in U$. Jest to zbiór
borelowski, bo warunek sprawdzamy dla
przeliczalnej liczby przyrostów. Z różniczkowalności
wynika $\bigcup_lH_{j,l}=D\cap Q$.
Dla każdego $j$ wybierzmy zatem $l_j$ tak duże,
aby $\lambda_n((D\cap Q)\setminus H_{j,l_j})
<\varepsilon2^{-j-2}$. Zbiór
$G=\bigcap_{j\ge1}H_{j,l_j}$ spełnia
$\lambda_n(Q\setminus G)<\varepsilon/2$.
Z twierdzenia Łuzina, zastosowanego do mierzalnej
funkcji macierzowej $DF$ na $G$, i z wewnętrznej
regularności miary Lebesgue'a otrzymujemy zwarty
$K\subseteq G$, dla którego
$\lambda_n(G\setminus K)<\varepsilon/2$ oraz
$DF|_K$ jest ciągła. Jeśli $K=\varnothing$,
przyjmujemy $F_K\equiv0$ i pierwsza teza jest
trywialna; w dalszym ciągu zakładamy $K\ne\varnothing$.

Warunek w definicji $H_{j,l_j}$, przez ciągłość
$F$ względem przyrostu $h$, pozostaje prawdziwy
również dla rzeczywistych $h$ o tej samej
ograniczonej długości. Stąd dla każdego $j$ i
$x,y\in K$ o $0<|x-y|<l_j^{-1}$ zachodzi
\[
 |F(y)-F(x)-DF(x)(y-x)|\le j^{-1}|x-y|.
\]
Zatem dane $a=F|_K$ i $A=DF|_K$ spełniają
założenia Lematu~\ref{lem:whitney-c1-28}.
Jego zastosowanie daje żądane przedłużenie.

Aby uzyskać ostatnią tezę, wybieramy przeliczalne
pokrycie $U$ zwartymi kostkami $Q_s\Subset U$.
Dla każdego $s$ i $j$ wykonujemy konstrukcję
z $\varepsilon=1/j$ na $Q_s$, uzyskując
$K_{s,j}$. Ponieważ
$\lambda_n(Q_s\setminus\bigcup_jK_{s,j})\le1/j$
dla każdego $j$, miara ta jest zerowa.
Przeliczalna suma pozostałych zbiorów zerowych
ma miarę zero, co kończy dowód.
\end{proof}

\section{Dwa oszacowania dla zbiorów krytycznych}
\begin{lemat}[Obraz zbioru o niedostatecznej randze]
\label{lem:critical-area-null-28}
Niech $1\le k\le N$, $V\subseteq\mathbb R^k$ będzie
otwarty, $G\in C^1(V;\mathbb R^N)$, a
$K\subseteq V$ zwarty. Jeżeli
$\operatorname{rank}DG(x)<k$ dla wszystkich $x\in K$,
to $\mathcal H^k(G(K))=0$.
\end{lemat}
\begin{proof}
Niech $M$ ogranicza normę $DG$ na pewnym
zwartym otoczeniu $K$. Z jednostajnej ciągłości
$DG$ wynika, że dla każdego $\eta>0$ i dostatecznie
małej kostki $Q$ o boku $r$, przecinającej $K$,
po wybraniu $x_Q\in K\cap Q$ mamy
\[
 |G(x)-G(x_Q)-DG(x_Q)(x-x_Q)|\le\eta\sqrt k\,r
 \quad(x\in K\cap Q).
\]
Obraz $G(K\cap Q)$ leży więc w otoczeniu o
grzubości $C_k\eta r$ zbioru zawartego
w afinicznej podprzestrzeni wymiaru co najwyżej
$k-1$ i promienia $C_kMr$. Taki zbiór można
pokryć najwyżej $C_{k,N,M}\eta^{-(k-1)}$
kulami promienia $C_k\eta r$ (siatka w owej
podprzestrzeni); suma ich $k$-tych potęg
średnic nie przekracza $C_{k,N,M}\eta r^k$.

Ponieważ $K$ mieści się w ustalonej ograniczonej
kostce, liczba kostek siatki o boku $r$
przecinających $K$ jest nie większa niż $Cr^{-k}$.
Otrzymujemy pokrycie $G(K)$ zbiorami o średnicy
nie większej niż $C\eta r$, którego suma
$k$-tych potęg średnic jest nie większa niż
$C_{k,N,M}\eta$. Gdy $r\downarrow0$, dopuszczalne
średnice także dążą do zera, więc
$\mathcal H^k(G(K))\le C_{k,N,M}\eta$.
Przechodząc $\eta\downarrow0$, dostajemy tezę.
\end{proof}

\begin{lemat}[Zbiory zerowe i ich włókna]
\label{lem:null-fibres-lipschitz-28}
Niech $n\ge m\ge1$, $F:\mathbb R^n\to\mathbb R^m$
będzie lokalnie Lipschitzowska, a
$Z\subseteq\mathbb R^n$ zbiorem o zerowej zewnętrznej
mierze Lebesgue'a. Wówczas
\[
 \mathcal H^{n-m}(Z\cap F^{-1}(\{y\}))=0
 \quad\text{dla }\lambda_m\text{-prawie wszystkich }y.
\]
\end{lemat}
\begin{proof}
Przez przeliczalne pokrycie dziedziny otwartymi
kostkami, na których $F$ ma wspólną stałą
Lipschitza $L$, wystarczy rozpatrzyć część $Z$
zawartą w jednej ograniczonej takiej kostce.
Piszemy $d=n-m$. Ustalmy $\delta,\varepsilon>0$.
Z warunku $\lambda_n^*(Z)=0$ wynika istnienie
przeliczalnego pokrycia $Z$ kostkami $Q_i$
o średnicach mniejszych niż $\delta$ i takiego, że
$\sum_i(\operatorname{diam}Q_i)^n<\varepsilon$.
Kostki możemy wybrać wewnątrz otoczenia,
na którym działa stała $L$.
Niech $z_i$ będzie środkiem $Q_i$ i niech
$V_i=B(F(z_i),L\operatorname{diam}Q_i)$.
Jeśli $y\in F(Z\cap Q_i)$, to $y\in V_i$.
Z definicji miary Hausdorffa z ograniczeniem
średnic uzyskujemy punktowo
\[
 \mathcal H^d_\delta(Z\cap F^{-1}(\{y\}))
 \le C_{n,m}\sum_i(\operatorname{diam}Q_i)^d\chi_{V_i}(y).
\]
Dla $d=0$ koszt niepustego zbioru w pokryciu
wynosi $1$. Prawa strona jest borelowska,
a całkując ją, otrzymujemy
\[
 \int_{\mathbb R^m}C_{n,m}\sum_i
 (\operatorname{diam}Q_i)^d\chi_{V_i}(y)
 \,\mathrm d\lambda_m(y)
 \le C_{n,m}L^m\varepsilon.
\]
Zatem dla dowolnego $t>0$ zbiór, na którym
$\mathcal H^d_\delta$ rozpatrywanego włókna
przekracza $t$, ma zewnętrzną miarę nie większą
niż $C_{n,m}L^m\varepsilon/t$. Ponieważ
$\varepsilon$ jest dowolne, miara ta jest zerowa.
Stosując to do $\delta=1/j$ i sumując
przeliczalnie zbiory wyjątkowe, dostajemy
$\mathcal H^d=0$ na prawie wszystkich włóknach.
\end{proof}

\begin{lemat}[Włókna zwartego zbioru krytycznego]
\label{lem:critical-coarea-null-28}
Niech $n\ge m\ge1$, $V\subseteq\mathbb R^n$
będzie otwarty, $G\in C^1(V;\mathbb R^m)$,
a $K\subseteq V$ zwarty. Jeżeli
$\operatorname{rank}DG(x)<m$ dla wszystkich $x\in K$,
to
\[
 \mathcal H^{n-m}(K\cap G^{-1}(\{y\}))=0
 \quad\text{dla }\lambda_m\text{-prawie wszystkich }y.
\]
\end{lemat}
\begin{proof}
Piszemy $d=n-m$. Najpierw zauważmy, że dla
$\delta>0$ funkcja
$y\mapsto\mathcal H^d_\delta(K\cap G^{-1}(\{y\}))$
jest borelowska. Rzeczywiście, dla $t>0$ warunek,
że jej wartość jest mniejsza niż $t$, można
wyrazić istnieniem skończonego pokrycia włókna
otwartymi zbiorami będącymi skończonymi
sumami kul o wymiernych środkach i promieniach,
średnicach mniejszych niż $\delta$ i sumie kosztów
mniejszej niż $t$. Dla ustalonego takiego pokrycia
$\mathcal U$ zbiór parametrów $y$, dla których
całe włókno mieści się w $\mathcal U$, jest
$\mathbb R^m\setminus G(K\setminus\mathcal U)$,
a więc jest otwarty, gdyż $K\setminus\mathcal U$
jest zwarty. Przeliczalna suma tych otwartych
zbiorów daje borelowskość. W konsekwencji
również funkcja dla $\delta\downarrow0$ jest
borelowska jako granica monotoniczna.

Ustalmy $\eta>0$. Pokryjmy $K$ kostkami siatki
o boku $r$, dla których wybieramy
$x_Q\in K\cap Q$. Gdy $r$ jest dostatecznie
mały, jednostajna ciągłość $DG$ daje
\[
 |G(x)-G(x_Q)-DG(x_Q)(x-x_Q)|\le C_n\eta r
 \quad(x\in K\cap Q).
\]
Ponieważ ranga $DG(x_Q)$ jest mniejsza niż $m$,
zbiór $G(K\cap Q)$ zawiera się w
$C_n\eta r$-otoczeniu afinicznej podprzestrzeni
wymiaru nie większego niż $m-1$, w obrębie
kuli o promieniu $C_nMr$, gdzie $M$ jest
lokalnym ograniczeniem $\|DG\|$.
Porównanie objętości po rozłożeniu
$\mathbb R^m$ na tę podprzestrzeń i jej
ortogonalne dopełnienie daje
\[
 \lambda_m(G(K\cap Q))\le C_{n,m,M}\eta r^m.
\]
Dla każdego $y$ zbiór $K\cap G^{-1}(\{y\})$
można pokryć fragmentami leżącymi w tych
kostkach $Q$, dla których $y\in G(K\cap Q)$.
Każdy fragment ma średnicę nie większą niż
$\sqrt n\,r$, więc dla $\delta_r=2\sqrt n\,r$
\[
 \mathcal H^d_{\delta_r}(K\cap G^{-1}(\{y\}))
 \le C_{n,m}\sum_Q r^d\chi_{G(K\cap Q)}(y).
\]
W sumie występuje najwyżej $Cr^{-n}$ kostek.
Całkując i stosując oszacowanie ich obrazów,
dostajemy
\[
 \int_{\mathbb R^m}\mathcal H^d_{\delta_r}
 (K\cap G^{-1}(\{y\}))\,\mathrm d\lambda_m(y)
 \le C_{n,m,M}\eta.
\]
Przechodząc $r\downarrow0$ (wzdłuż malejącego
ciągu) i korzystając z monotoniczności
$\mathcal H^d_\delta$ oraz twierdzenia o
zbieżności monotonicznej, otrzymujemy
$\int\mathcal H^d(K\cap G^{-1}(\{y\}))
\,\mathrm d\lambda_m(y)\le C\eta$.
Dowolność $\eta$ kończy dowód.
\end{proof}


\section{Wzory area i coarea dla odwzorowań Lipschitzowskich}
W poniższych wzorach $DF$ oznacza różniczkę z twierdzenia
Rademachera, istniejącą prawie wszędzie. Jakobian $J_kF$
oraz jakobian normalny $\operatorname{Jac}_mF$ obliczamy
w punktach różniczkowalności, podstawiając $DF(x)$
do funkcji jakobianowych z
Definicji~\ref{def:k-jacobian-28} i
Definicji~\ref{def:normal-jacobian-28},
a na pozostałym zbiorze zerowym przypisujemy im wartość $0$.
Tak określone funkcje są mierzalne, a wybór wartości
na zbiorze wyjątkowym nie zmienia całek.
\begin{twierdzenie}[Wzory area i coarea dla odwzorowań Lipschitzowskich]
\label{thm:lipschitz-area-coarea-28}
Niech $E\subseteq\mathbb R^k$ będzie mierzalny
w sensie Lebesgue'a, a $F:\mathbb R^k\to\mathbb R^N$,
$k\le N$, będzie lokalnie Lipschitzowski.
Wówczas różniczka $DF(x)$ istnieje dla prawie
wszystkich $x$ (twierdzenie Rademachera)
i dla każdej mierzalnej $f\ge0$ zachodzi
\[
 \int_E f(x)J_kF(x)\,\mathrm d\lambda_k(x)
 =\int_{\mathbb R^N}\left(\sum_{x\in E\cap F^{-1}(\{y\})}
 f(x)\right)\mathrm d\mathcal H^k(y).
\]
Ponadto jeżeli $n\ge m\ge1$ i
$F:\mathbb R^n\to\mathbb R^m$ jest lokalnie
Lipschitzowski, to dla każdej mierzalnej
funkcji $g\ge0$ zachodzi
\[
 \int_{\mathbb R^n}g(x)\operatorname{Jac}_mF(x)
 \,\mathrm d\lambda_n(x)
 =\int_{\mathbb R^m}\left(
 \int_{F^{-1}(\{y\})}g(x)
 \,\mathrm d\mathcal H^{n-m}(x)\right)
 \,\mathrm d\lambda_m(y).
\]
W przypadku $n=m$ wewnętrzna całka po włóknie
jest sumą względem miary liczącej.
\end{twierdzenie}

\begin{lemat}[Wzór area dla dowolnego odwzorowania $C^1$]
\label{lem:c1-general-area-28}
Niech $1\le k\le N$, $V\subseteq\mathbb R^k$ będzie
otwarty, $G\in C^1(V;\mathbb R^N)$, a $A\subseteq V$
borelowski. Dla każdej nieujemnej funkcji borelowskiej
$f:A\to[0,+\infty]$ zachodzi
\[
 \int_Af(x)J_kG(x)\,\mathrm d\lambda_k(x)
 =\int_{\mathbb R^N}\left(
 \sum_{x\in A\cap G^{-1}(\{y\})}f(x)\right)
 \,\mathrm d\mathcal H^k(y).
\]
\end{lemat}
\begin{proof}
Zbiór $V_{\mathrm{reg}}=\{x\in V:\operatorname{rank}DG(x)=k\}$
jest otwarty, a $G|_{V_{\mathrm{reg}}}$ jest immersją.
Stosujemy do niego
Twierdzenie~\ref{thm:immersion-area-28}.
Jego dopełnienie w $V$ jest względnie domkniętym
zbiorem punktów krytycznych, który można wyczerpać
przeliczalną rodziną zwartych podzbiorów.
Lemat~\ref{lem:critical-area-null-28} dowodzi, że obraz
tych punktów jest $\mathcal H^k$-zerowy.
Na części krytycznej $J_kG=0$, a suma po włóknach
może być różna od zera tylko na wspomnianym
zbiorze zerowym. Równości po zsumowaniu obu
części są więc prawdziwe.
\end{proof}

\begin{proof}[Dowód ogólnych wzorów area i coarea z
Twierdzenia~\ref{thm:lipschitz-area-coarea-28}]
\emph{Krok 1: wzór area.}
Z Lematu~\ref{lem:lusin-whitney-28} otrzymujemy
zwarte zbiory $K_j\subseteq\mathbb R^k$ oraz
odwzorowania $F_j\in C^1(\mathbb R^k;\mathbb R^N)$
takie, że $F_j=F$ i $DF_j=DF$ na $K_j$, a
$Z=\mathbb R^k\setminus\bigcup_jK_j$ jest
$\lambda_k$-zerowy. Połóżmy
$A_1=K_1$ i
$A_j=K_j\setminus\bigcup_{i<j}K_i$ dla $j\ge2$.
Są to parami rozłączne zbiory borelowskie,
których suma jest dopełnieniem $Z$.

Dla $E$ i $f$ borelowskich stosujemy
Lemat~\ref{lem:c1-general-area-28} do
$F_j$ na $A_j\cap E$. Daje on
\[
 \int_{A_j\cap E}f(x)J_kF(x)\,\mathrm d\lambda_k(x)
 =\int_{\mathbb R^N}
 \sum_{x\in A_j\cap E\cap F^{-1}(\{y\})}f(x)
 \,\mathrm d\mathcal H^k(y).
\]
Po zsumowaniu względem $j$ stosujemy Tonellego.
Zbiór $Z$ nie wnosi nic do lewej strony.
Ponieważ $F$ jest lokalnie Lipschitzowskie,
$\mathcal H^k(F(Z))=0$: dla każdego zwartego
fragmentu dziedziny wynika to z nierówności
$\mathcal H^k(F(A))\le L^k\mathcal H^k(A)$
z Rozdziału~\ref{ch:hausdorff-26}.
Zatem również prawa strona nie zmienia się
po dołączeniu punktów z $Z$.

Jeżeli $E$ jest tylko mierzalny w sensie
Lebesgue'a, a $f$ jest mierzalna, wybieramy
borelowskie przedstawiciele $E$ i $f$ różniące
się od zadanych jedynie na zbiorze
$\lambda_k$-zerowym. Jego obraz przez $F$
jest $\mathcal H^k$-zerowy, więc ta modyfikacja
nie zmienia żadnej strony równości. Zapewnia
to również mierzalność prawej strony względem
uzupełnionej miary $\mathcal H^k$.

\emph{Krok 2: wzór coarea na częściach pełnej rangi.}
Przyjmijmy na razie, że $g$ jest borelowska.
Dla odwzorowania $F:\mathbb R^n\to\mathbb R^m$
wybieramy analogiczne $K_j$ i $F_j\in C^1$,
z rozłącznymi zbiorami $A_j\subseteq K_j$.
Niech
$V_j=\{x:\operatorname{rank}DF_j(x)=m\}$.
Jest to zbiór otwarty, a $F_j|_{V_j}$ jest
submersją. Twierdzenie~\ref{thm:coarea-submersion-28}
zastosowane do funkcji
$g\chi_{A_j\cap V_j}$ daje
\begin{align*}
 &\int_{A_j\cap V_j}g(x)\operatorname{Jac}_mF(x)
 \,\mathrm d\lambda_n(x)\\
 &\quad=\int_{\mathbb R^m}\left(
 \int_{A_j\cap V_j\cap F^{-1}(\{y\})}
 g(x)\,\mathrm d\mathcal H^{n-m}(x)\right)
 \,\mathrm d\lambda_m(y).
\end{align*}

\emph{Krok 3: części krytyczne i zbiór wyjątkowy.}
Zbiór
$C_j=K_j\cap\{x:\operatorname{rank}DF_j(x)<m\}$
jest zwarty. Na nim $\operatorname{Jac}_mF=0$,
a Lemat~\ref{lem:critical-coarea-null-28}
pokazuje, że
$\mathcal H^{n-m}(C_j\cap F^{-1}(\{y\}))=0$
dla prawie wszystkich $y$. Na pozostałym zbiorze
$Z=\mathbb R^n\setminus\bigcup_jK_j$ jakobian
również nie wnosi nic do całki, bo
$\lambda_n(Z)=0$. Z kolei
Lemat~\ref{lem:null-fibres-lipschitz-28} dowodzi,
że $Z$ jest $\mathcal H^{n-m}$-zerowy na
prawie wszystkich włóknach. Stąd po zsumowaniu
równości z kroku 2 otrzymujemy pełny wzór coarea.

W przypadku funkcji $g$ jedynie mierzalnej
w sensie Lebesgue'a zastępujemy ją wersją
borelowską poza zbiorem zerowym. Lemat o
zerowych włóknach gwarantuje, że taka zmiana
nie wpływa na prawą stronę dla prawie wszystkich
$y$; jej mierzalność wynika z rozpatrzenia
części borelowskich i uzupełnienia miary.
Nie trzeba w żadnym kroku zakładać, że odwzorowanie
$F$ jest globalnie różnowartościowe lub że ma
stałą rangę.
\end{proof}

\begin{uwaga}[Zakres uzyskanego dowodu]
W porównaniu z wersjami gładkimi nowym elementem
jest jednoczesna kontrola zbiorów krytycznych i ich
włókien. Lemat~\ref{lem:whitney-c1-28} wyprowadza
potrzebne przedłużenie $C^1$, a lematy
\ref{lem:critical-area-null-28}--\ref{lem:critical-coarea-null-28}
wyprowadzają oszacowania bez powoływania się na
same wzory area/coarea. Dzięki temu dowód nie jest
błędnym kołem. Nadal poza zakresem pozostają
analogiczne wzory na ogólnych przestrzeniach
metrycznych oraz dla odwzorowań, które nie są
lokalnie Lipschitzowskie.
\end{uwaga}

\section{Postać na rozmaitościach}
Niech $(M^n,g)$ i $(N^m,h)$ będą rozmaitościami
riemannowskimi, $n\ge m$, a $F:M\to N$
submersją $C^1$. Różniczkę i jej sprzężenie
$\mathrm dF_p^*$ wyznaczają metryki $g,h$.
Jakobian normalny definiujemy jako
$\operatorname{Jac}_{g,h}F(p)=
\sqrt{\det(\mathrm dF_p\mathrm dF_p^*)}$,
gdzie wyznacznik liczymy w bazie ortonormalnej
$T_{F(p)}N$. Dla nieujemnej borelowskiej $\varphi$
analogiczny lokalny rachunek daje
\begin{equation}\label{eq:coarea-manifolds-28}
 \int_M\varphi(p)\operatorname{Jac}_{g,h}F(p)
 \,\mathrm d\mu_g(p)
 =\int_N\left(\int_{F^{-1}(\{q\})}\varphi(p)
 \,\mathrm d\mu_{g|_{F^{-1}(\{q\})}}(p)\right)
 \,\mathrm d\mu_h(q).
\end{equation}
Dowód redukuje się w mapach do
Twierdzenia~\ref{thm:coarea-submersion-28}:
zmiana objętości w mapach źródła, celu i włókna
prowadzi do tej samej tożsamości wyznaczników
co w Lemacie~\ref{lem:linear-coarea-28}, zapisanej
w bazach ortonormalnych. Lokalne równości
łączymy przeliczalnym podziałem borelowskim.
Nie wykorzystujemy orientowalności.

\begin{wniosek}[Miara riemannowska i Hausdorffa]
\label{cor:riemann-hausdorff-28}
Na gładkiej rozmaitości riemannowskiej $(M^n,g)$
miara $\mu_g$ jest równa $n$-wymiarowej mierze
Hausdorffa zbudowanej z odległości riemannowskiej,
przy konwencji kalibracji z
Uwagi~\ref{rem:hausdorff-convention-26}. Na osadzonej
gładkiej podrozmaitości $S^k\subseteq M$ zachodzi
analogiczna identyfikacja miary indukowanej
$\mu_{g|_S}$ z $\mathcal H^k$ obliczaną
z odległości ograniczonej z $M$ na $S$.
\end{wniosek}
\begin{proof}
Wokół każdego punktu wybieramy mapę normalną
lub dowolną mapę, w której metryka ma ciągłe
współczynniki. Na dostatecznie małym otoczeniu
metryczne odległości są, z błędem mnożnikowym
$1+\varepsilon$, porównywalne z odległościami
określonymi przez stały iloczyn skalarny
$g_p$. Wynika to z ciągłości $g$ i
lokalnego porównania długości krzywych
(zmniejszając otoczenie tak, aby minimalizujące
krzywe nie opuszczały większego otoczenia).
Miara Hausdorffa tych samych zbiorów różni się
wtedy co najwyżej o czynnik $(1+\varepsilon)^n$
od miary Hausdorffa dla metryki stałej.
Ta ostatnia ma gęstość $\sqrt{\det G(p)}$
względem $\lambda_n$. Przez lokalną
regularność miar i przejście $\varepsilon\downarrow0$
otrzymujemy równość miar w mapie. Przeliczalne
pokrycie mapami daje ją globalnie.
Dla $S$ powtarzamy argument dla metryki
indukowanej; odległość wewnętrzna na $S$
i zewnętrzna odległość ograniczona z $M$
są lokalnie asymptotycznie równoważne,
więc ich miary Hausdorffa wymiaru $k$ są równe.
\end{proof}

\begin{uwaga}[Co przechodzi do geometrycznej teorii miary?]
Następny rozdział zastosuje miary Hausdorffa,
formuły area i coarea oraz twierdzenia o
różniczkowaniu miar do zbiorów prostowalnych
i zbiorów o skończonym obwodzie. Te obiekty
nie muszą być gładkimi podrozmaitościami,
a zatem stanowią właściwe rozszerzenie
geometrii różniczkowej.
\end{uwaga}

\begin{uwaga}[Dżety przybliżone a wzory miarowo-geometryczne]
Pochodna $DF(x)$ z twierdzenia Rademachera jest częścią
pierwszego dżetu odwzorowania $F$ w punktach różniczkowalności.
Przy osłabieniu pojęcia granicy do granicy przybliżonej otrzymujemy
\emph{dżety przybliżone}, określone miarowo prawie wszędzie.
W rozdziale~\ref{ch:dzety} rozwijamy geometryczną teorię wiązek
dżetów, omawiamy twierdzenie Whitneya o przedłużaniu danych
dżetowych oraz jego miarową wersję Whitneya--Luzina.
Jest to uzupełnienie obecnego rozdziału, a nie ponowny dowód
wzorów area i coarea.
\end{uwaga}
